\section{Lean formalization architecture}
\label{sec:lean}

The formalization plan is tied to the actual repository declarations rather
than an invented surrogate calculus.  This section gives the main definitions,
exact theorem endpoints, and dependency order needed for a complete Lean
implementation.

\subsection{Existing interfaces}

The intended imports are:

\begin{lstlisting}[language={},caption={Repository interfaces used by the compiler}]
import ConciseIPRA.Syntax
import ConciseIPRA.Deduction
import ConciseIPRA.PR
\end{lstlisting}

The namespace is \texttt{LeanProofs.ConciseIPRA}.  The relevant types and
constructors are:

\begin{lstlisting}[language={}]
PR k                         -- primitive-recursive symbol
PTm n                        -- term with n free variables
Fm PTm n                     -- first-order formula
PRAAx                        -- iPRA axioms
Prf PRAAx Gamma phi          -- embedded derivability
PRA.qfInd                    -- quantifier-free induction
Prf.substWeaken, Prf.trans   -- substitution/weakening and cut
PRA.precZero, PRA.precSucc   -- recursion equations
\end{lstlisting}

The checked packet module \path{ObjectIteration.lean} should be integrated as
\path{ConciseIPRA.FiniteTrace}; its key declarations are
\path{primitiveRecursionInvariant} and \path{primitiveRecursionAt}.  The
checked \path{ObjectObservation.lean} combinators should similarly become a
library module \path{ConciseIPRA.ObservableGlue}.  Integration should import
the repository sources directly rather than retain frozen copies of the
kernel.

\subsection{Trace package}

The accompanying \texttt{ProofCarryingGraph.lean} proposes the following
structure.

\begin{lstlisting}[language={}]
structure TracePackage (k : Nat) (Gamma : List (Fm PTm k)) where
  init        : PR k
  step        : PR (k + 2)
  invariant   : Fm PTm (k + 2)
  invariantQF : invariant.QF
  base        : Gamma |-PRA baseFormula init invariant
  advance     : invariant :: (Gamma.map shift).map shift
                  |-PRA nextFormula step invariant
\end{lstlisting}

Here \texttt{|-PRA} abbreviates the repository notation
\(\Gamma\vdash[\PRAAx]\varphi\).  The method
\texttt{TracePackage.runAt} is a direct wrapper around
\path{primitiveRecursionAt}; it contains no new mathematics.  A
\texttt{TraceExit} stores fuel, output formula, and an object proof from the
terminal invariant to the output.  Its compiler is one simultaneous cut.

This layer should be the first integration milestone because it immediately
gives future application modules a standard type for bounded verified runs.

\subsection{Proof-producing graph rules}

A graph rule can be represented without dependent graph syntax.

\begin{lstlisting}[language={}]
structure Rule (n : Nat) where
  inputs : List (Fm PTm n)
  output : Fm PTm n
  derive : forall Gamma,
    (forall A, A `mem` inputs -> Gamma |-PRA A) ->
    Gamma |-PRA output
\end{lstlisting}

A compiled prefix stores a list of formulas and a proof of every member.
Extending it with a rule is a single host-level construction.  This design has
several advantages:

\begin{itemize}
\item the graph compiler is polymorphic in the ambient context;
\item dependencies are ordinary formula membership, not a custom proof
language;
\item rules can close over arbitrary theorem-specific data and existing
proofs;
\item the final result is visibly an ordinary \texttt{Prf PRAAx} term.
\end{itemize}

A more strongly typed graph DSL may be added later, but it is not needed for
the central theorem.

\subsection{Term-to-symbol compilation}

The graph compiler must emit \texttt{PR} symbols, not arbitrary host
functions.  The resources-and-elimination packet reports checked declarations
\path{TermCompiler.symbol_correct} and
\path{TermCompiler.transport_equation}.  These should be integrated and
used to compile fixed term expressions to symbols with object equation proofs.

The compiler should support:
\begin{enumerate}
\item variables and constants;
\item composition;
\item pairing, projections, and fixed finite tuples;
\item primitive recursion;
\item bounded conditionals and Boolean operations;
\item finite table reads and writes;
\item fixed reverse-index transformations.
\end{enumerate}

A host function returning \texttt{PR k} is acceptable because it constructs
syntax.  A host function \(\N^k\to\N\) is not automatically acceptable; it
must be accompanied by compilation and equations.

\subsection{Boolean reflection packages}

The blueprint defines

\begin{lstlisting}[language={}]
structure BoolReflection (phi : Fm PTm n) where
  arity    : Nat
  code     : PR arity
  args     : Fin arity -> PTm n
  forward  : [phi] |-PRA checked code args
  backward : [checked code args] |-PRA phi
\end{lstlisting}

The implementation order is:

\begin{enumerate}
\item prove equality and Boolean operation equations;
\item define \texttt{allBelow}, \texttt{anyBelow}, and \texttt{leastBelow};
\item prove their object reflection by \texttt{PRA.qfInd};
\item define a host inductive syntax \texttt{BForm};
\item recursively compile every constructor to a \texttt{BoolReflection}.
\end{enumerate}

Implication needs decidability of the recursively compiled subformulas, not
classical excluded middle.  The reflection fields provide the required case
split.

\subsection{Finite table library}

A reusable table module should expose:

\begin{lstlisting}[language={}]
Table.empty     : PR ...
Table.read      : PR ...
Table.write     : PR ...
Table.length    : PR ...
Table.mapFold   : PR ...
Table.all       : PR ...
Table.any       : PR ...
\end{lstlisting}

and object derivations for read-after-write, read-other, length, fold equations,
and bounded validity.  The exact encoding can be pairing-based or sequence-
number based.  What matters is a small equation interface.

For fixed graph size, an even simpler tuple code suffices.  The general table
library is useful for variable finite lists such as point enumerations and
candidate clouds.

\subsection{Demand compiler}

Define a host structure for a fixed node specification:

\begin{lstlisting}[language={}]
structure NodeSpec where
  predecessors : Array Nat
  demand       : Array (PR demandArity)
  localRule    : Rule parameterArity
  -- plus syntax for evaluator/certificate fields
\end{lstlisting}

The host compiler checks topological order, constructs the reverse-pass PR
symbol, and emits formulas for each table row.  The object proof has two
levels.

\begin{enumerate}
\item Defining equations of the reverse recurrence, reusing the checked
\path{DemandCompiler.object_backward_zero} and
\path{DemandCompiler.object_backward_succ} results where their interfaces match.
\item A quantifier-free invariant stating that the completed suffix of the
demand table satisfies all edge inequalities.
\end{enumerate}

The demand leastness theorem can be a separate object theorem.  Forward
correctness needs only soundness; leastness is valuable for complexity and
review.

\subsection{Finite dispatch}

A trace step for a fixed graph must dispatch on a program counter.  The
recommended route is a compiled selector term and one finite theorem:

\begin{quote}
If \(pc<N\) and branch \(i\) preserves the invariant for every numeral
\(i<N\), then the selector-dispatched step preserves the invariant.
\end{quote}

Because \(N\) is fixed by the host graph, the proof constructor can recursively
nest Boolean equality cases and apply the corresponding branch derivation.  It
does not need induction on arbitrary program syntax.

Alternatively, generate a specialized step proof directly as a nested
\texttt{orE} tree.  This may produce larger terms but is simple and robust.

\subsection{Relational viability module}

The viability recurrence should use reverse time \(r=H-i\), avoiding
subtraction in the primitive-recursion argument.  Define a PR table
\(W(H,r,a,\vec x)\) whose row \(r\) corresponds to level \(H-r\).  The step
uses \texttt{anyBelow} over the next cloud.

Main targets:

\begin{lstlisting}[language={}]
viable_iff_path
leastViable_start
leastViable_edge
leastViable_valid
leastPath_lexicographic
shrinking_step_bound
shrinking_cauchy
\end{lstlisting}

The path predicate is itself a bounded checker over a table code.  Thus
\texttt{viable\_iff\_path} is a quantifier-free induction theorem after
Booleanization.

\subsection{Rational metric layer}

Use normalized integer pairs or signed numerators with positive denominators.
Required exact operations:

\begin{itemize}
\item equality and order by cross multiplication;
\item absolute difference;
\item powers of two;
\item finite sums and maxima;
\item interval containment and width;
\item triangle inequality and telescoping.
\end{itemize}

The shrinking theorem needs no completed real-number type.  It proves a
first-order inequality between rational observations.  A ``regular name'' is
just a fixed PR symbol returning rational codes plus the embedded Cauchy
formula.

\subsection{End-to-end theorem shape}

Every application should expose its final endpoint in a small audit module.
For example:

\begin{lstlisting}[language={}]
def cubicObservableSentence : Fm PTm 0 := ...

theorem cubicObservable :
    ([] : List (Fm PTm 0)) |-PRA cubicObservableSentence := by
  exact (compileGraph cubicGraph).head ...

#print axioms cubicObservable
\end{lstlisting}

The audit should reject weaker endpoints such as a host theorem
\texttt{CubicObservableSpec}, an uninstantiated goal structure, or a theorem
that assumes the desired \texttt{Prf} term.

\subsection{Recommended module graph}

\begin{center}
\begin{tabular}{@{}ll@{}}
\toprule
Module & Responsibility \\
\midrule
\path{ConciseIPRA.FiniteTrace} & integrated checked trace rule \\
\path{ConciseIPRA.ObservableGlue} & integrated checked certificate combinators \\
\texttt{PCOG.TermCompiler} & fixed term to \texttt{PR} plus equations \\
\texttt{PCOG.Boolean} & Boolean operations and reflection \\
\texttt{PCOG.BoundedFold} & all/any/least bounded search \\
\texttt{PCOG.Table} & finite coded tables and invariants \\
\texttt{PCOG.Demand} & reverse demand pass \\
\texttt{PCOG.Graph} & \texttt{Rule}, \texttt{Compiled}, graph fold \\
\texttt{PCOG.Cloud} & viability and least finite paths \\
\texttt{PCOG.RationalMetric} & intervals, error algebra, telescoping \\
\texttt{PCOG.Stable} & independent witnesses and gap collapse \\
\texttt{PCOG.Examples.Cubic} & first end-to-end benchmark \\
\bottomrule
\end{tabular}
\end{center}

\subsection{Formalization risks}

The principal risks are engineering rather than conceptual.

\begin{itemize}
\item Scoped de Bruijn substitutions can make invariant instantiation verbose.
The existing trace module already solves the hardest state/time substitution
pattern; reuse it.
\item Large nested conjunctions can cause proof-term and simplifier blowup.
Use Boolean table checkers instead of formula-level conjunctions where
possible.
\item Host arrays and functions are not automatically \texttt{PR}.  Keep a
clear boundary between compile-time graph traversal and emitted syntax.
\item A theorem that proves a semantic evaluator correct does not supply its
object equations.  Require equation derivations in every compiler result.
\item Full normalization or transport to the native scalar core should be a
final independent step via the repository's existing equivalence; do not mix
it into the PCOG compiler.
\end{itemize}

\subsection{Verification protocol}

A completed implementation should record:

\begin{enumerate}
\item repository commit and source hashes;
\item Lean version and root Lake build output;
\item \texttt{\#print axioms} for every public embedded theorem;
\item absence of \texttt{sorry}, new opaque axioms, and
\texttt{Classical.choice} in audited dependencies;
\item theorem types confirming \texttt{Prf PRAAx}, not host analogues;
\item finite executable tests only as supplementary evidence;
\item a theorem ledger marking each target as definition, human proof,
compiled host theorem, embedded proof, or conditional interface.
\end{enumerate}

The files in this archive follow the ledger discipline, but no Lean build is
claimed for them in the present environment.
